%*****************************************************************************
% Layout-Vorschlag zur Erstellung von Bachelor-, Master oder Diplomarbeiten  *
% Modifikation moeglich                                                      *
% S. Gippert, 2012                                                           *
% ****************************************************************************

\documentclass[11pt,a4paper]{scrbook}
\usepackage{graphicx}
% \usepackage{ngerman}
\usepackage[ngerman]{babel}
\usepackage{amsmath}
\usepackage{amssymb}
\usepackage{amsthm}

\usepackage{multirow}
\usepackage{booktabs}

\usepackage{hyperref}

% Umgebungen
\renewcommand{\baselinestretch}{1.50}\normalsize
\newenvironment{Beweis}{\begin{proof}[Beweis]}{\end{proof}}
\newtheorem{Satz}{Satz}[chapter]% Definitionen
\def\Real{{{\rm I\!R}}}

% Dokumentbeginn
\begin{document}

% Titelseite ----------ANPASSEN!----------------------
\begin{titlepage}
\begin{center}
\huge \textbf{{Layer-Wise Relevance Propagation For Deep Neural Networks}} \\
\vspace{2cm}
\LARGE\textbf{{Abschlussbericht}}\\
\vspace{2cm}
\normalsize


vorgelegt von  \textbf{Gruppe 4\\ Theo Conrads, Robin Kühling, Marc Bremser}\\
am \today \\
\vspace{2.5cm}
\textbf{am Mathematischen Institut der \\Universit\"at zu K\"oln}\\
\vspace{3cm}
\end{center}
\normalsize{
\begin{center}
Erstgutachter:  {Prof.~Dr.~Axel Klawonn} \\
\vfill
\includegraphics[scale=0.25]{uzk_logo} % bunt
% \includegraphics[scale=0.25]{uzk_logo_grau} % grau
\end{center}
}
\end{titlepage}


\pagestyle{plain}
\cleardoublepage
\setcounter{page}{1}
\pagenumbering{roman}
\tableofcontents

\listoffigures
\addcontentsline{toc}{chapter}{Abbildungsverzeichnis}

\listoftables
\addcontentsline{toc}{chapter}{Tabellenverzeichnis}

\cleardoublepage

% Dateien
\chapter{Einleitung}
Hier grundsätzliche Einführung
\section{Taylor Decomposition}
\frame{\sectionpage}
\frame{\frametitle{Taylor Decomposition - Rückblick}
\begin{itemize}
\item Einfaches Netzwerk mit einem Hidden Layer, ReLU Aktivierung und Sum-Pooling als Output.
\item Zusätzliche Voraussetzung: $b_j \leq 0$. 
\begin{figure}[H]
\includegraphics[width = 0.7 \textwidth]{grafiken_marc/simple_net.png}
\end{figure}
\item Für das Outputneuron $x_k$ gilt: $x_k = max(0,\sum_{j} x_j)$
%\item Gesucht ist eine Vorschrift, nach der die Relevanz $R_k$ auf die Neuronen der vorherigen Schicht verteilt wird.
\end{itemize}
}

\frame{\frametitle{Taylor Decomposition - Rückblick}
\begin{itemize}
\item Suche eine Nullstelle f\"ur die Taylorentwicklung von $R_k(\boldsymbol{x}) = \sum_j x_j$.

\begin{figure}
\includegraphics[width = 0.6\textwidth]{grafiken_marc/rel_prop_1.png}
\end{figure}
\pause
\item Wg. ReLU Aktivierung im vorherigen Layer und $\sum_j x_j \overset{!}{=} 0$ ist $\tilde{\boldsymbol{x}}=$\textbf{0} die einzige Nullstelle von $R_k$.
\item Wegen $R_j=\frac{\partial R_{k}}{\partial x_{j}}(x_j - \tilde{x}_j)=1 \cdot (x_j - 0)$ gilt also
\item $R_{j}=x_{j}= \max \left(0, \sum_{i} x_{i} w_{i j}+b_{j}\right)$
\end{itemize}
}

\frame{\frametitle{Taylor Decomposition - Generische Regel}
\begin{itemize}
%\item Führe nun für jede der $R_j$ Taylorentwicklung mit einem eigenen Entwicklungspunkt $\hat{x}^{(j)}$ durch.
\item Es gilt $R_j = x_j = \max \left(0, \sum_{i} x_{i} w_{i j}+b_{j}\right)$

\begin{figure}
\includegraphics[width = 0.6\textwidth]{grafiken_marc/rel_prop_2.png}
\end{figure}
%\pause
\item Unterscheide nun 2 Fälle:
\begin{enumerate}
\item $R_j=0$: Nicht aktivierte Neuronen sollen keine Relevanz zurückverteilen. Insbesondere gilt hier $\tilde{\bx}= \bx $.
\item $R_j>0$: Hierfür wird ein Richtungsvektor $\bv^{(j)}$ definiert. $\tilde{\bx}$ soll von der Form $\bx	+ t \cdot \bv^{(j)}$, mit $t \in \rn$ sein.
\end{enumerate}
\end{itemize}
}


\begin{frame}{Taylor - Entwicklungspunkt}
\begin{itemize}
\item Allgemeine Vorgehensweise: 
\item Durch Einsetzen von $\tilde{\bx}=\bx	+ t \cdot \bv^{(j)}$ in die Ebenengleichung $\sum_{i} \tilde{x}_{i} w_{i j}+b_{j} = 0$ lässt sich eine allgemeine Formel für $t$ finden.
\item Somit gilt: 
\begin{align*}
0 &= \sum_{i} \left(x_{i} + t v_i^{(j)} \right) w_{i j}+b_{j} \\
\Leftrightarrow -t& =\frac{\sum_{i} x_{i} w_{i j}+b_{j}}{\sum_{i} v_{i}^{(j)} w_{i j}}\\
\ \\
\Rightarrow x_i - \tilde{x_i} &= -t v_i^{(j)} = \frac{\sum_{i} x_{i} w_{i j}+b_{j}}{\sum_{i} v_{i}^{(j)} w_{i j}} v_i^{(j)}
\end{align*}
\end{itemize}
\end{frame}

\begin{frame}{Taylor - Entwicklungspunkt}
\begin{itemize}
\item Für die Umverteilung von der $l+1$-ten Schicht in die $l$-te Schicht gilt
\begin{align*}
R_i^l &= \sum_{j} R_{i \leftarrow j}^l =  \sum_{j} \frac{\partial R_{j}^{l+1}}{\partial x_{i}^l}
(x_i - \tilde{x_i}) \\
&= \sum_{j : R_j = 0} \frac{\partial R_{j}^{l+1}}{\partial x_{i}^l} \cdot 0 + \sum_{j: R_j>0} w_{ij} \frac{\sum_{i} x_{i} w_{i j}+b_{j}}{\sum_{i} v_{i}^{(j)} w_{i j}} v_i^{(j)} \\
&= \sum_{j} \frac{v_i^{(j)}  w_{i j}}{\sum_{i} v_{i}^{(j)} w_{i j}} R_j^{l+1}
\end{align*}
\item $\Rightarrow$ Allgemeine Formel in Abhängigkeit von $v_i^{(j)}$
\end{itemize}
\end{frame}


\section{Min-Max Modell}
Min-Max-Modell?
%\section{Part 2}
\begin{frame}{Erweiterung auf tiefe Netze}
\begin{itemize}
\item Bei tiefen Netzen, insbesondere Convolutional Layern ist die Relevanzfunktion nicht unbedingt explizit angegeben.
\end{itemize}
\vspace*{20pt}
\begin{minipage}{0.42\textwidth}
\includegraphics[width=\textwidth]{grafiken_marc/simple_net.png}
\end{minipage}
\begin{minipage}{0.52\textwidth}
\includegraphics[width=\textwidth]{grafiken_marc/DNN.png}
\end{minipage}
\vspace*{10pt}
\begin{itemize}
\item Ein Feature kann vorhanden sein, aber bei der Bilderkennung keine Rolle spielen
\end{itemize}
\footnote{Grafik entnommen aus \ref{itm:Mont17}}
\end{frame}

\begin{frame}{Erweiterung auf tiefe Netze}
\begin{itemize}
\item Gesucht ist eine Approximation der Relevanz Funktion, die leicht zu analysieren ist.
\item Führe das Konzept \textbf{Relevanz-Modell} ein, um bei tieferen Netzen die Relevanzfunktion $R_j^{l+1}(\bx^{l})$ zu approximieren.
\item Nehme an, die Relevanzfunktion $R_j$ lässt sich in der Form $R_j = x_j \cdot c_j$ schreiben, wobei $c_j$ konstant ist.
\item Im Paper als "{}Training Free"{} Ansatz vorgestellt
\end{itemize}
\end{frame}

\begin{frame}{Erweiterung auf tiefe Netze}
\begin{itemize}
\item Nehme an, die Relevanzfunktion $R_j$ lässt sich in der Form $R_j = x_j \cdot c_j$ schreiben, mit $c_j$ konstant.
\item Betrachte die generische Redistributionsregel
\begin{align*}
R_i &= \sum_{j} \frac{x_{i} \cdot \rho(w_{i j})}{\sum_{i} x_{i} \cdot \rho(w_{i j})} R_{j}\\
& = x_{i} \sum_{j}  \frac{ \rho(w_{i j})}{\sum_{i} x_{i} \cdot \rho(w_{i j})} x_j \cdot c_j \\
&= x_i \sum_{j}   \rho(w_{i j}) \frac{\max \left(0, \sum_{i} x_{i} w_{i j}\right)}{\sum_{i} x_{i} \cdot \rho(w_{i j})} c_j
\end{align*}
\end{itemize}
\end{frame}

\begin{frame}{Erweiterung auf tiefe Netze}
\begin{itemize}
\item Nehme an, die Relevanzfunktion $R_j$ lässt sich in der Form $R_j = x_j \cdot c_j$ schreiben, mit $c_j$ konstant.
\item Betrachte die generische Redistributionsregel
\begin{align*}
R_i &= \sum_{j} \frac{x_{i} \cdot \rho(w_{i j})}{\sum_{i} x_{i} \cdot \rho(w_{i j})} R_{j}\\
& = x_{i} \sum_{j}  \frac{ \rho(w_{i j})}{\sum_{i} x_{i} \cdot \rho(w_{i j})} x_j \cdot c_j \\
&= x_i \underbrace{\sum_{j}   \rho(w_{i j}) \frac{\max \left(0, \sum_{i} x_{i} w_{i j}\right)}{\sum_{i} x_{i} \cdot \rho(w_{i j})} c_j}_{\colonapprox c_i}
\end{align*}
\end{itemize}
\end{frame}

\begin{frame}{Erweiterung auf tiefe Netze}
\begin{itemize}
\item Es gilt also
\begin{align*}
R_i = x_i \underbrace{\sum_{j}   \rho(w_{i j}) \frac{\max \left(0, \sum_{i} x_{i} w_{i j}\right)}{\sum_{i} x_{i} \cdot \rho(w_{i j})} c_j}_{\colonapprox c_i} = \sum_{j}    \frac{\rho(w_{i j}) \cdot x_i}{\sum_{i} x_{i} \cdot \rho(w_{i j})} R_j
\end{align*}
\item D.h. $R_i^l$ lässt sich wieder schreiben als $x_i^l \cdot c_i^l$, mit $c_i^l$ annähernd konstant.
\item Ausgehend von der letzten Schicht kann die Relevanz somit auch gemäß der hergeleiteten Regeln zurück zum Input verteilt werden.
\item Der Parameter $c_i^l$ wird dabei durch die betrachtete Regel "{}induktiv"{} gebildet.
\end{itemize}
\end{frame}


\begin{frame}{Zusammenhang mit LRP}
\begin{itemize}
\item Die klassische $LRP-0$ Formel kann als Deep Taylor Entwicklung im Nullpunkt gesehen werden
\item Betrachte O.B.d.A. ein Neuron $x_j$ mit $R_j > 0$. 
\item Die Suchrichtung $\bv$ ist hierbei der Punkt $\bx$ selbst, und somit gilt:
\pause
\begin{align*}
R_{i \leftarrow j}^l &=\frac{\partial R_{j}^{l+1}}{\partial x_{i}^l}
(x_i - \tilde{x_i})  = w_{ij} \cdot c_j \cdot (x_i - \tilde{x_i}) \\
&= w_{ij} \cdot c_j \cdot \frac{\sum_{i} x_{i} w_{i j}+b_{j}}{\sum_{i} x_{i} w_{i j}} x_i\\
&= \frac{x_i \cdot w_{ij}}{\sum_{i} x_{i} w_{i j}} x_j \cdot c_j
\end{align*}
\end{itemize}
\end{frame}


\begin{frame}{Zusammenhang mit LRP}
\begin{itemize}
\item Die klassische $LRP-0$ Formel kann als Deep Taylor Entwicklung im Nullpunkt gesehen werden
\item Betrachte O.B.d.A. ein Neuron $x_j^{l+1}$ mit $R_j > 0$. 
\item Die Suchrichtung $\bv$ ist hierbei der Punkt $\bx$ selbst, und somit gilt:
\begin{align*}
R_{i \leftarrow j}^l &=\frac{\partial R_{j}^{l+1}}{\partial x_{i}^l}
(x_i - \tilde{x_i})  = w_{ij} \cdot c_j \cdot (x_i - \tilde{x_i}) \\
&= w_{ij} \cdot c_j \cdot \frac{\sum_{i} x_{i} w_{i j}+b_{j}}{\sum_{i} x_{i} w_{i j}} x_i\\
&= \frac{x_i \cdot w_{ij}}{\sum_{i} x_{i} w_{i j}} \underbrace{ x_j \cdot c_j}_{ R_j }
\end{align*}
\end{itemize}
\end{frame}

\begin{frame}{Zusammenhang mit LRP}
\begin{itemize}
\item Für die totale Relevanz von $x_i$ gilt:

\begin{align*}
R_i^l = \sum_{j} R_{i \leftarrow j}^l &= \sum_{j} \frac{x_i \cdot w_{ij}}{\sum_{i} x_{i} w_{i j}} R_j^{l+1} = \sum_{j} \frac{z_{ij}}{\sum_{i} z_{i j}} R_j^{l+1}
\end{align*}
\item Mit der gleichen Vorgehensweise können auch die $LRP-\varepsilon$ Regel sowie $LRP-\gamma$ hergeleitet werden.
\item Deep Taylor Decomposition als Basis für \textit{LRP}
\item Wesentlicher Unterschied: Geforderte Konsistenz von Montavon et al. im Vergleich zu den \textit{LRP} Regeln
\end{itemize}
\end{frame}


\section{Implementierung - Netzarchitektur}
\begin{itemize}
\item Netwerkarchitekturen, den Programmcode vorstellen
\end{itemize}
%\section{Part 2}
\begin{frame}{Aufgabe}
\begin{itemize}
\item Herleitung $z^B$ Regel und Implementierung der einzelnen Regeln
\item Vergleich der Ergebnisse
\end{itemize}
\end{frame}
\section{Min-Max Implementierung}
\begin{itemize}
\item Implementierung Min-Max Modell
\end{itemize}

% Literaturverzeichnis
\bibliography{thesisbib}{}
\bibliographystyle{plain}
\addcontentsline{toc}{chapter}{Literaturverzeichnis}
%%\newpage
\thispagestyle{empty}
\chapter*{Erkl\"arung}
\addcontentsline{toc}{chapter}{Erkl\"arung}
\noindent
Hiermit versichere ich an Eides statt, dass ich die vorliegende Arbeit selbstst\"andig und ohne die Benutzung anderer als der angegebenen Hilfsmittel angefertigt habe. Alle Stellen, die w\"ortlich oder sinngem\"a\ss   aus ver\"offentlichten und nicht ver\"offentlichten Schriften entnommen wurden, sind als solche kenntlich gemacht. Die Arbeit ist in gleicher oder \"ahnlicher Form oder auszugsweise im Rahmen einer anderen Pr\"ufung noch nicht vorgelegt worden. Ich versichere, dass die eingereichte elektronische Fassung der eingereichten Druckfassung vollst\"andig entspricht.
\begin{tabular}{p{5cm} p{7cm}}
\vspace{4cm}\\
\hline
Ort, Datum & Unterschrift
\end{tabular}


\end{document}
